\documentclass[../../main.tex]{subfiles}

\begin{document}

\subsection{microFrameWork} \label{subsec:il-micro-framework}
  microFrameWork è un framework \textit{object-oriented} minimale scritto
  interamente in \CPP{}
  che permette di creare dei \textit{workflow} di analisi dati scrivendo
  delle configurazioni, sotto forma di file json, che descrivono come
  concatenare moduli di input, output e elaborazione, implementati
  anch'essi in \CPP.

  Dati un insieme di \texttt{TTree} con lo stesso numero di entries contenuti
  in uno o più file .root, il ufw è progettato per leggere da tutti ed
  elaborare un indice alla volta, per poi scrivere nello stesso indice di
  altri tree in uno o più file di output.
  Ogni indice può essere processato in maniera indipendente dagli altri, ma
  nonostante ciò attualmente il framework non permette di
  elaborare più indici in parallelo.\\
  Il termine tecnico per questi indici è \texttt{context\_id}.\\
  Il ufw prevede anche la possibilità di leggere oggetti che non siano
  \texttt{TTree}, ad esempio il \texttt{TGeoManager} citato nella
  sezione~\ref{sec:edep-sim}, in questi casi sarà usato lo stesso oggetto
  per ogni indice.

  La natura \textit{object-oriented} del ufw è ravvisabile nel fatto che
  ogni elemento che deve essere gestito da esso necessita di ereditare da
  una classe virtuale definita nel ufw stesso.
  Queste classi sono tre, una per ogni tipo di oggetto necessario al
  framework: \texttt{ufw::process}, \texttt{ufw::data::base} e
  \texttt{ufw::streamer}.

  \subsubsection*{I processi}
    I processi rappresentano gli step di elaborazione dati del
    \textit{workflow} di analisi.
    Quanto questi step debbano essere ``grandi'' è a discrezione
    dell'utente ma la prassi vorrebbe che si trovi il giusto mezzo tra
    l'avere un \textit{workflow} composto da un solo modulo e avere dei
    moduli che eseguono poche righe di codice ognuno.

    I processi possono accettare sia parametri di configurazione indicati
    dall'utente nel json, sia dati provenienti da altri processi o
    direttamente da file (approfondimento su questo punto in seguito) e
    possono a loro volta ritornare dei dati in output.\\
    Per creare un processo bisogna creare una classe che eredita da
    \texttt{ufw::process}.

  \subsubsection*{I dati}
    Nel contesto del ufw con il termine dati ci si riferisce a tutti e i
    soli oggetti ai quali possono accedere i vari moduli e streamer, la cui
    \texttt{ownership} appartiene al framework.
    Non rientrano sotto questa definizione, per esempio, tutte le variabili
    create all'interno dei singoli moduli, che non devono categoricamente
    essere passate ad altri oggetti del framework.

    Per rappresentare un tipo di dato, una classe deve ereditare dalla
    classe template \texttt{ufw::data::base<T1, T2, T3>}, dove i tre
    parametri template, detti tag, servono a specificare alcune
    caratteristiche del dato.

    \paragraph{Tag 1: Managed vs Complex}
      Questo tag specifica le modalità di I/O del tipo.\\
      Per i dati contrassegnati come Managed l'I/O è gestito da uno
      streamer, mentre per quelli Complex è richiesto all'utente che
      li dichiara di creare anche una specializzazione del template
      \texttt{ufw::data::factory<T>}
      (dove \texttt{T} è la classe che si sta implementando) che sulla
      base della configurazione indicata nel json andrà a costruire
      un'istanza del tipo in questione assumendone la \texttt{ownership}.
      Non è possibile usare i tipi dichiarati Complex come output del
      \textit{workflow}.

    \paragraph{Tag 2: Unique vs Instanced}
      Questo tag specifica quante istanze del tipo in questione possono
      essere create.\\
      Come suggerito dal nome, un tipo contrassegnato come Unique ammette
      una sola istanza, alla quale non è assegnato alcun nome, per farvi
      riferimento nel file di configurazione basta indicare il tipo stesso,
      i tipi Instanced al contrario ammettono più istanze, ognuna delle
      quali, nel json, è identificata solo da un nome (non dal tipo).\\
      Al momento della scrittura ci sono due precisazioni importanti da
      tenere a mente:
    \begin{itemize}
      \item l'unicità dei tipi Unique non è garantita in alcun modo dal
      framework;
      \item i controlli sulla corrispondenza dei tipi richiesti dai moduli
      e quelli forniti dal framework (sulla base della configurazione
      json) sono effettuati solo a \textit{runtime} tramite dei
      \texttt{dynamic\_cast}.
      \end{itemize}

    \paragraph{Tag 3: Global vs Context}
      Questo tag specifica il \textit{life time} delle istanze di questo
      tipo.\\
      Le istanze dei tipi contrassegnati come Global vengono costruite al primo
      \texttt{context\_id} e distrutte al termine dell'elaborazione
      dell'ultimo, per quanto riguarda i tipi Context, come suggerisce il
      nome, il loro \textit{life time} è legato a quello del
      \texttt{context\_id}.\\
      Questo tag è quello usato per distinguere i dati letti (scritti) dai
      (nei) due tipi di oggetti che il ufw è in grado di gestire, discussi prima
      (\texttt{TTree} e non), ma non è esclusivo dei dati di input o output
      (per quanto ciò sarà comune).

  \subsubsection*{Gli streamer}
    Gli streamer sono responsabili delle interfacce di input e output del
    framework.
    Per la precisione il ufw può creare o modificare un file in output solo
    tramite uno streamer, mentre per quanto riguarda l'input c'è anche la
    possibilità di usare i sopracitati dati contrassegnati dal tag Complex.

    L'utente che dichiara uno streamer, oltre a implementare un metodo
    \texttt{read} e/o un metodo \texttt{write} in base al tipo di interfaccia
    che si vuole creare, è anche responsabile della creazione degli eventuali
    dizionari ROOT o affini, necessari alle operazioni di I/O dei tipi che
    intende gestire.\\
    Per rappresentare uno streamer, una classe deve ereditare da
    \texttt{ufw::streamer}.


\end{document}